\newcommand{\auditoriaid}{A05 -- Arquitectura y run.sh}
% Preambulo comun de las auditorias del informe E1 (revision_e1/)
\documentclass[11pt,a4paper]{article}
\usepackage[utf8]{inputenc}
\usepackage[T1]{fontenc}
\usepackage[spanish,es-noshorthands]{babel}
\usepackage{lmodern}
\usepackage{microtype}
% Sin patrones de guionado en espanol en esta instalacion (ver A10): se
% desactiva el guionado automatico para evitar cortes con reglas del ingles.
\hyphenpenalty=10000
\exhyphenpenalty=10000
\usepackage[a4paper,margin=2.2cm]{geometry}
\usepackage{booktabs,array,longtable}
\usepackage{graphicx,caption,subcaption}
\usepackage{xcolor}
\usepackage{enumitem}
\usepackage{fancyhdr}
\usepackage{amsmath}
\usepackage{tikz}
\usetikzlibrary{arrows.meta,positioning}
\usepackage{natbib}
\usepackage{xurl}
\usepackage[hidelinks]{hyperref}
\urlstyle{tt}
\newcommand{\archivo}[1]{\url{#1}}
\newcolumntype{L}[1]{>{\raggedright\arraybackslash}p{#1}}
\definecolor{alta}{RGB}{180,30,30}
\definecolor{media}{RGB}{200,120,0}
\definecolor{baja}{RGB}{60,110,160}
\definecolor{cajafondo}{RGB}{245,247,250}
\newcommand{\sevA}{\textcolor{alta}{\textbf{Alta}}}
\newcommand{\sevM}{\textcolor{media}{\textbf{Media}}}
\newcommand{\sevB}{\textcolor{baja}{\textbf{Baja}}}
\setlength{\parskip}{0.35em}
\setlength{\emergencystretch}{3em}
\setlength{\parindent}{0pt}
\pagestyle{fancy}
\fancyhf{}
\fancyhead[L]{\small Auditoría del Informe del Entregable~1 -- \auditoriaid}
\fancyhead[R]{\small \thepage}
\renewcommand{\headrulewidth}{0.4pt}
% Entorno de hallazgos: ID | Severidad | Ubicacion | Hallazgo y evidencia | Correccion
\newenvironment{hallazgos}{%
\small
\begin{longtable}{@{}L{1.15cm}L{1.25cm}L{1.9cm}L{6.0cm}L{\dimexpr\textwidth-10.3cm-10.6\tabcolsep\relax}@{}}
\toprule
\textbf{ID} & \textbf{Sev.} & \textbf{Ubicación} & \textbf{Hallazgo y evidencia} & \textbf{Corrección aplicada} \\
\midrule\endfirsthead
\toprule
\textbf{ID} & \textbf{Sev.} & \textbf{Ubicación} & \textbf{Hallazgo y evidencia} & \textbf{Corrección aplicada} \\
\midrule\endhead
\bottomrule\endfoot
}{\end{longtable}}
% Macros compartidas con el informe revisado (las secciones propuestas las usan)
% Macros compartidas entre el informe revisado y las auditorias.
% Cifras verificadas contra los archivos del repositorio (ver A01/A04/A07):
\newcommand{\nUniverso}{102}   % source_id CMIP6 con tos/Omon en ESGF (00b, model_availability_priority.csv)
\newcommand{\nSeed}{41}        % publican los cuatro experimentos (config/models_seed_cmip6.csv)
\newcommand{\nSel}{40}         % completos + QC PASS (models_inventory_final.csv)
\newcommand{\nComb}{160}       % combinaciones modelo-experimento PASS (qc_report.csv)
\newcommand{\nNoEnc}{44}       % config/models_no_encontrado_44.csv
\newcommand{\nParcial}{18}     % 102 - 40 - 44
% Fecha de la portada: fija (no \today), editable en un solo lugar.
\newcommand{\fechaentrega}{28 de septiembre de 2026}
% Estado de codigo que documenta el E1 (antes de los cambios del E2)
\newcommand{\commitEone}{\texttt{8ee6184}}

% En las auditorias el texto propuesto se muestra fuera del informe: las
% referencias cruzadas se imprimen con el nombre de la seccion destino.
\makeatletter
\def\etq#1#2{\expandafter\def\csname etq@#1\endcsname{#2}}
\etq{sec:resumen}{Resumen}
\etq{sec:alcance}{Alcance}
\etq{sec:marco}{Revisión científica}
\etq{sec:antecedentes}{Antecedentes}
\etq{sec:referencias}{Literatura sobre TOE}
\etq{tab:referencias}{bibliografía TOE}
\etq{sec:similitudes}{Niño~3.4 frente a Niño~1+2}
\etq{sec:toe_metodos}{Elección de los métodos de TOE}
\etq{tab:metodos_toe}{métodos TOE}
\etq{eq:szabo_signal}{ecuación señal M1}
\etq{eq:szabo_V}{ecuación V(t)}
\etq{eq:szabo_T}{ecuación T(t)}
\etq{eq:szabo_toe}{ecuación TOE M1}
\etq{eq:tod}{ecuación ToD}
\etq{sec:datos}{Datos}
\etq{fig:cmip6_contribution}{contribución institucional}
\etq{sec:tabla_modelos}{Modelos seleccionados}
\etq{tab:modelos}{modelos}
\etq{sec:desviacion}{Selección de modelos}
\etq{sec:dominios_periodos}{Dominios y periodos}
\etq{fig:region_nino}{regiones Niño}
\etq{sec:arquitectura}{Arquitectura del pipeline}
\etq{fig:flujograma}{Flujograma}
\etq{sec:scripts}{Descripción de los pasos}
\etq{sec:paso00b}{Paso 00b}
\etq{sec:paso01}{Paso 01}
\etq{sec:paso02}{Paso 02}
\etq{sec:paso04}{Paso 04}
\etq{sec:paso05}{Paso 05}
\etq{sec:paso06}{Paso 06}
\etq{sec:paso07}{Paso 07}
\etq{sec:reporte}{Reporte de estado}
\etq{sec:resultados}{Resultados}
\etq{fig:qc}{matriz de estado}
\etq{tab:cuentas}{conteo de modelos}
\etq{fig:mapas2d}{mapas 2D}
\etq{fig:series}{series}
\etq{fig:boxplot}{boxplots}
\etq{sec:roadmap}{Próximos pasos}
\etq{tab:roadmap}{próximos pasos}
\etq{sec:conclusiones}{Conclusiones}
\etq{sec:recomendaciones}{Recomendaciones}
\makeatother

\makeatletter
\AtBeginDocument{\renewcommand{\ref}[1]{\@ifundefined{etq@#1}{\textit{[#1]}}{\textit{\csname etq@#1\endcsname}}}}
\makeatother
% Texto propuesto: secciones degradadas un nivel y en caja
\newenvironment{propuesto}{%
  \par\medskip\noindent\textcolor{baja}{\rule{\textwidth}{0.8pt}}\par
  \noindent{\small\textbf{Inicio del texto propuesto para el informe revisado}}\par\medskip
  \begingroup
  \let\section\subsection\let\subsection\subsubsection\let\subsubsection\paragraph
}{%
  \endgroup
  \par\noindent{\small\textbf{Fin del texto propuesto}}\par
  \noindent\textcolor{baja}{\rule{\textwidth}{0.8pt}}\par\medskip}
% Recuadro de actualizacion posterior (decisiones del autor)
\newcommand{\actualizacion}[1]{\par\medskip\noindent\fcolorbox{media}{media!8}{\parbox{\dimexpr\linewidth-2\fboxsep-2\fboxrule\relax}{\small\textbf{Actualización del 29-09-2026 (decisiones del autor).} #1}}\par\medskip}

\begin{document}

\begin{center}
{\Large\bfseries Auditoría A05: Arquitectura del \emph{pipeline} y \texttt{run.sh}}\\[0.4em]
{\normalsize Informe del Entregable~1, Sección~5 (\texttt{informe\_e1\_smnv3.tex}, commit \texttt{330fd2e})}\\[0.2em]
{\small Revisión: \fechaentrega}
\end{center}

\section{Objeto}

Se audita la Sección~5 (``\emph{Pipeline}'': principios, ``Script principal \texttt{run.sh}'' y flujograma). Por instrucción del autor, la información de \emph{uso} del proyecto (instalación, argumentos, ejemplos, tiempos, salidas) sale del informe y pasa a un documento de anexos separado, escrito para un usuario básico de Linux. Esta auditoría define qué queda en el informe, qué pasa al anexo y qué debe corregirse. Todo se contrastó con \texttt{run.sh}, \texttt{environment.yml} y los scripts auxiliares en el commit \texttt{8ee6184}.

\actualizacion{El flujograma pasó a la Sección~6 (``Obtención y procesamiento de los datos''), que ahora es un párrafo con las acciones realizadas y el flujograma explicado brevemente. La Sección~5 conserva los principios de diseño, sin remitir a pasos individuales.}

\section{Comportamiento verificado de \texttt{run.sh} (commit \texttt{8ee6184})}

\begin{enumerate}[leftmargin=1.4em]
\item Activa por sí mismo el entorno conda \texttt{smn\_toe} (o \texttt{e1\_smn}, el nombre anterior) si existe, y antepone su \texttt{bin} al \texttt{PATH}, por los clústeres HPC donde el \texttt{python3} del sistema tiene prioridad.
\item Valida los argumentos \texttt{[STEP\_FROM] [STEP\_TO] [MAX\_MODELS]}; los pasos válidos son \texttt{00 00b 01 02 03 04 05 06 07}.
\item Muestra un \textbf{reporte \emph{preflight}} (\texttt{preflight\_report.py}): qué artefactos existen, cuáles faltan y cuánto tardaron en corridas anteriores en esa máquina. El informe original no lo menciona.
\item Ejecuta los pasos del rango pedido, registra la salida en \texttt{logs/pipeline.log} y la duración de cada paso en \texttt{logs/step\_timings.csv}.
\item Siempre, sin importar el rango: registro de modelos, cinco scripts de gráficos, reporte de estado (\archivo{logs/run_report_<fecha>.txt}) y paquete \archivo{informe/entregable_update.tar.gz}.
\end{enumerate}

Hay dos preguntas interactivas, ambas con límite de 15\,s. La de buscar en nodos alternativos y Copernicus la hace \texttt{02\_download\_all\_sources.sh}, no \texttt{run.sh}; solo aparece con terminal interactiva y modelos pendientes, y el silencio equivale a ``sí''. La de regenerar las figuras existentes la hace \texttt{plot\_common.py}; el silencio equivale a ``no''.

\textbf{Recursos medidos en esta máquina} (\texttt{logs/step\_timings.csv}, \texttt{logs/download\_timings.csv}): Paso~00b, mediana de 95\,min; Paso~01, hasta 45\,min; Paso~02 (descarga de 160 combinaciones), hasta $\sim$13\,h en una corrida, con 164 descargas medidas, mediana de 3,5\,min y máximo de 84\,min cada una; Paso~04, hasta 49\,min; Pasos~05--07, pocos minutos. Espacio: \texttt{data/raw/cmip6} 73\,GB, \texttt{data/interim} 4,7\,GB y \texttt{data/processed/masked} 4,3\,GB (esta última medida con la ventana ampliada del E2).

\section{Hallazgos}

\begin{hallazgos}
A05-1 & \sevM & Títulos 5 y 6 & Dos secciones consecutivas con el mismo título, ``\emph{Pipeline}''. & ``Arquitectura del \emph{pipeline}'' (5) y ``Descripción de los pasos'' (6). \\
A05-2 & \sevM & 5.1 & El detalle de uso de \texttt{run.sh} (argumentos, ejemplos, guardia de descargas, reportes) ocupa unas 600 palabras en el cuerpo del informe. & Por instrucción del autor pasa al Anexo; el informe conserva un párrafo sobre el papel de \texttt{run.sh} y remite al anexo. \\
A05-3 & \sevM & 5 (principios) & ``un \texttt{git clone} nuevo más \texttt{./run.sh}, sin intervención manual, obtiene todo lo que sea posible''. Omite dos requisitos previos (crear el entorno conda; credenciales de Copernicus para 02c) y que desde \texttt{main} se obtiene el estado del E2, no el del E1 (ver A01-6). & En el informe la frase se reemplaza por los tres principios; los requisitos y la versión que se debe usar se explican en el Anexo. \\
A05-4 & \sevM & 5 (principios) & La idempotencia se presenta como absoluta, pero el Paso~06 la abandona a propósito desde \texttt{c07f4c5}. & Se nombra la excepción y su razón. \\
A05-5 & \sevM & 5.1 & No menciona el reporte \emph{preflight}, la activación automática del entorno ni el registro de tiempos, que son lo primero que ve un usuario al ejecutar \texttt{run.sh}. & Se documentan en el Anexo. \\
A05-6 & \sevB & Flujograma & El código TikZ conserva nodos comentados (``Notebook'', ``Reporte + paquete final'') y un comentario erróneo (\texttt{scale=0.95 \% reduce todo al 75\%}). El pie menciona el \emph{notebook} de gráficos exploratorios, que dejó de versionarse en \texttt{9ba3f2c} y que el diagrama no dibuja. & Código limpio; la última caja pasa a ser ``Gráficos, reporte y paquete''; el pie se ajusta a lo dibujado. \\
A05-7 & \sevB & Código (fuera del informe) & El comentario de cabecera de \texttt{run.sh} está desactualizado: nombra LLNL como nodo principal de 00b (ahora es ORNL), \texttt{model\_availability\_report.*} como salida de 00b (ahora es \texttt{model\_availability\_priority.csv}) y el \emph{notebook} como vía de gráficos. & No se modifica código en esta revisión. Se recomienda actualizarlo junto con el cierre del E1 (ver A10). \\
A05-8 & -- & 5.1 & Verificado sin cambios: orden fijo de pasos; \texttt{logs/pipeline.log}; invocación del Paso~02 a través de \texttt{run\_download\_if\_idle.sh} con \texttt{pgrep}; cinco scripts \texttt{plot\_*.py}; reporte de estado y paquete siempre al final. & Se conservan (en el informe o en el Anexo). \\
\end{hallazgos}

\section{Reparto de contenido entre informe y anexo}

\begin{center}\small
\begin{tabular}{@{}L{7.4cm}L{\dimexpr\textwidth-7.4cm-2\tabcolsep\relax}@{}}
\toprule
Queda en el informe (Sección~5) & Pasa al Anexo de uso \\
\midrule
Repositorio, punto de entrada, salidas en \texttt{data/processed/}; tres principios de diseño; papel de \texttt{run.sh} en un párrafo; flujograma. & Requisitos (Linux, conda, disco, red, credenciales CDS); instalación; argumentos y ejemplos; qué muestra el \emph{preflight}; preguntas interactivas; tiempos; carpetas de salida; cómo leer el reporte de estado; problemas frecuentes; cómo reproducir el estado del E1. \\
\bottomrule
\end{tabular}
\end{center}

\section{Extensión}

Original: 749 palabras (incluido el flujograma). Propuesto para el informe: 307 
palabras. El resto se traslada al Anexo, donde se amplía con los detalles que necesita un usuario básico.

\section{Texto propuesto}

\begin{propuesto}
% ============================================================
\section{Arquitectura del \emph{pipeline}}
\label{sec:arquitectura}
% ============================================================

El \emph{pipeline} vive en el repositorio \archivo{smn_toe}, tiene un único punto de entrada (\texttt{run.sh}) y deja sus salidas en \texttt{data/processed/}, de donde el Entregable~2 las lee sin necesidad de conocer el detalle interno de este módulo. Se rige por tres principios:

\begin{itemize}
\item \textbf{Idempotencia.} Cada paso omite el procesamiento cuya salida ya existe, de modo que una corrida puede interrumpirse y reanudarse sin recalcular lo hecho. La excepción deliberada es el control de calidad, que se recalcula completo en cada corrida porque es barato y un resultado en caché puede quedar obsoleto.
\item \textbf{Verificación por valor y no por metadato.} Las unidades, los valores de relleno y la lista de meses se comprueban sobre los datos numéricos y no sobre los atributos declarados en los archivos. Existen archivos cuyo atributo contradice su contenido; por ejemplo, archivos que declaran grados Celsius y contienen valores en Kelvin.
\item \textbf{Cascada de fuentes con degradación explícita.} Si el nodo principal de ESGF no resuelve un modelo, se prueban automáticamente nodos alternativos y Copernicus CDS. Lo que ninguna fuente resuelve queda registrado en un reporte de estado, en lugar de fallar en silencio.
\end{itemize}

\texttt{run.sh} no contiene lógica de procesamiento propia: ejecuta los pasos 00, 00b, 01, 02, 03, 04, 05, 06 y 07 en ese orden fijo, o solo un rango de ellos, y registra la salida de cada uno en \texttt{logs/pipeline.log}. Al terminar cualquier corrida, completa o parcial, genera siempre el registro de modelos, las figuras automáticas, el reporte de estado y un paquete comprimido con los artefactos para actualizar el informe. Así, cada ejecución deja constancia verificable de qué modelos quedaron descargados, procesados o pendientes. La instalación, los argumentos de \texttt{run.sh}, los tiempos y el espacio en disco requeridos, y la interpretación de las salidas se describen en el documento de anexos que acompaña a este informe.

\end{propuesto}

\end{document}
